% TOP_PROBLEM: {"id": 1100116, "title": "Questions in Geometric Group Theory — Q 1.16", "category_id": 17, "category": {"id": 17, "name": "group_theory", "display_name": "Group Theory", "description": "Problems about groups, group actions, representations, and related algebraic structures.", "slug": "group-theory", "order_index": 17, "created_at": "2026-07-31T15:26:25.671Z"}, "status": "open", "problem_number": "AMR-010-0116", "set_id": 13, "statusQualification": "Excludes the trivial group, which immediately refutes the literal unrestricted statement; the source discussion explicitly treats nontrivial finite groups."}
% FIRST_POSTED: 2026-10-01
% ENTRY_KIND: statement_only
% UnsolvedMath ID 1100116; source catalogue code AMR-010-0116
% !TeX program = lualatex
% Definition and sources only; this file is not an LLM solution attempt.
\documentclass[11pt,a4paper]{article}
\usepackage[margin=25mm]{geometry}
\usepackage{fontspec}
\setmainfont{lmroman10-regular.otf}[BoldFont=lmroman10-bold.otf,
 ItalicFont=lmroman10-italic.otf,BoldItalicFont=lmroman10-bolditalic.otf]
\usepackage{amsmath,amssymb,mathtools}
\usepackage{unicode-math}
\setmathfont{latinmodern-math.otf}
\AtBeginDocument{\let\setminus\smallsetminus}
\usepackage{xurl,hyperref}
\hypersetup{colorlinks=true,urlcolor=blue}
\setcounter{secnumdepth}{0}
\setlength{\parindent}{0pt}
\setlength{\parskip}{5pt}
\setlength{\emergencystretch}{3em}
\begin{document}
\section*{Questions in Geometric Group Theory — Q 1.16}\label{1100116}
{\small UnsolvedMath ID 1100116 · AMR-010-0116 · Group Theory · AMR Open Problem Lists}
\subsection{Definitions and mathematical statement}
For a finitely generated group $H$, let $d(H)$ be the minimum cardinality of a generating set. Generators here generate as an abstract group, without taking any topological closure. For a fixed group $G$ and an integer $n\geq1$, let $G^n$ be its direct product with coordinatewise multiplication.

Assume that $G$ is nontrivial and word-hyperbolic: a Cayley graph for a finite generating set has uniformly thin geodesic triangles. The question asks whether
\[
 \lim_{n\to\infty}d(G^n)=\infty.
\]
Explicitly, for every integer $R\geq0$, must there exist $N=N(G,R)$ such that $d(G^n)>R$ for every $n\geq N$? The group remains fixed while the number of factors grows.

The projection $G^{n+1}\to G^n$ shows that these ranks form a nondecreasing sequence. Thus the assertion is equivalent to their being unbounded, and a counterexample would have one finite bound for all powers. No particular rate of divergence or bound uniform over hyperbolic groups is requested. The trivial group must be excluded: all its direct powers have rank zero. This exclusion is also consistent with the source's accompanying discussion of nontrivial finite groups. The problem is about independent direct-product coordinates, not the ranks of finite-index subgroups or of successive quotients of $G$.
\subsection{Short English statement}
For a fixed nontrivial hyperbolic group, must the number of generators of its direct powers grow without bound?
\subsection{Sources}
\begin{itemize}
\item[S1] ULAM.AI and the UnsolvedMath Contributors, supplied problems.json, record 1100116 (AMR-010-0116), AMR Open Problem Lists. Catalogue identity and source transcription; CC BY 4.0.
\url{https://huggingface.co/datasets/ulamai/UnsolvedMath}.
\item[S2] Bestvina - Questions in Geometric Group Theory (pdf) (2004), Question 1.16 (PDF page 4). Original source indexed by AMR-010-0116.
\url{https://www.math.utah.edu/~bestvina/eprints/questions-updated.pdf}.
\end{itemize}
\end{document}
